\documentclass[11pt,a4paper]{article}
\usepackage{operationsform}
\begin{document}
\OPStitle{質問設計チェックリスト}{TPL-CHK-022}{0.2.0}
\OPSsection{1}{確認項目}
\OPScheck{1}{質問が取材目的と科学的主張・根拠に結びついている}
\OPScheck{2}{方法・対象・数値・推論・仮説・将来構想を区分している}
\OPScheck{3}{開放型質問と確認型質問、追質問が設計されている}
\OPScheck{4}{確定していない前提を事実として押しつけていない}
\OPScheck{5}{研究の制約や相反する解釈を確認する問いを含む}
\OPScheck{6}{発言者の負担、答えない選択、社会的文脈に配慮している}
\OPScheck{7}{限られた時間に対して質問の優先順位と代替案がある}
\OPScheck{8}{質問改訂の担当と関連する出典IDが記録されている}
\OPSsection{2}{点検結果の記録}
\OPSblankshort{対象案件・対象版}
\OPSblankshort{担当者と点検日}
\OPSblanklong{未解決事項と対応}
\par\Needspace{4\baselineskip}\noindent\textbf{確認に用いた資料ID}\par\vspace{1mm}
{\fontsize{8.5}{13}\selectfont
\begin{tabularx}{\linewidth}{@{}XXXX@{}}
\textbf{対象・項目} & \textbf{内容・参照} & \textbf{担当・状態} & \textbf{備考}\\[2mm]
\strut & \strut & \strut & \strut\\[5mm]
\strut & \strut & \strut & \strut\\[5mm]
\end{tabularx}
}\vspace{2mm}\par
\end{document}
